\section{Normal Subgroups, Quotients, and Homomorphisms}

Let \(H \subgroup G\), when are two cosets of \(H\) in \(G\) equal?

If \(gH = g' H\), then \(g \in g' H\), so we can write \(g = g' h\) for some \(h \in H\). Hence, \(\para{g'}\inverse g \in H\). The argument is reversible.

Therefore, we have
\[
    gH = g' H \iff \para{g'}\inverse g \in H.
\]

We denote the set of \emph{left-cosets} of \(H \subgroup G\) by
\[
    G / H = \set{gH}{g \in G},
\]
and such set \emph{always exists} -- the normality of \(H\) in \(G\) gives this set a well-defined algebraic structure.

The next question is, if \(G / H\) carry a \emph{natural} group structure? A natural way to define a group operation is such that
\[
    g_1 H \cdot g_2 H = \para{g_1 g_2} H.
\]

But note this formula may depend on the way in which we write the cosets \(g_1 H\) and \(g_2 H\), i.e. the choice of the particular elements \(g_1, g_2 \in G\).

If \(g_2 H = g_2' H\), then by previously, \(g_2' = g_2 h\) for some \(h \in H\). Therefore, we have
\[
    g_1 H \cdot g_2' H = \para{g_1 g_2' H} = g_1 g_2 h H = g_1 g_2 H,
\]
so this is fine. However, if \(g_1 H = g_1' H\), then \(g_1' = g_1 h\) for some \(h \in H\). We need
\[
    g_1 g_2 H = g_1 h g_2 H.
\]

But this is not necessarily true for all subgroups \(H \subgroup G\) and \(g_1, g_2 \in G\). It is certainly true for abelian groups -- but there is a weaker equivalent condition for this hold to hold, i.e.,
\[
    g_2 \inverse h g_2 \in H,
\]
for all \(h \in H\) and \(g_2 \in G\), in order for the proposed rule to define a well-defined group structure.

\begin{definition}[Normality]
    \label{def:normality}%
    A subgroup \(H \subgroup G\) is \emph{normal}, if for all \(g \in G\), and \(h \in H\), we have \(g h g \inverse \in H\). We denote this as \(H \normalsub G\).
\end{definition}

\begin{definition}[Quotient Group]
    \label{def:quotient-group}%
    When \(H \normalsub G\), then \(G / H\) does inherit a group structure by the rule above. We refer to this as the \emph{quotient group of \(G\) by \(H\)}, where the identity is \(H\).
\end{definition}

We recall the definition of a homomorphism in group theory.

\begin{definition}[Group Homomorphism]
    \label{def:group-homomorphism}%
    Let \(G\) and \(H\) be groups. A \emph{homomorphism} from \(G\) to \(H\) is a function \(\functiondomain{f}{G}{H}\), such that for all \(g_1, g_2 \in G\),
    \[
        f\para{g_1 g_2} = f\para{g_1} f\para{g_2}.
    \]
\end{definition}

It follows that we must have \(f\para{e_G} = e_H\) by calculating \(f\para{e_G \cdot e_G} = f\para{e_G}\).

\begin{lemma}[Group Homomorphisms Respect Inverses]
    \label{lem:group-homomorphisms-inverses}%
    If \(\functiondomain{f}{G}{H}\) is a homomorphism, then \(f\para{g\inverse} = f\para{g}\inverse\).
\end{lemma}

\begin{proof}
    We find \(f\para{g g\inverse}\) in two ways. First, note
    \[
        f\para{g g\inverse} = f\para{e_G} = e_H.
    \]

    On the other hand, we have
    \[
        f\para{g g\inverse} = f\para{g} f\para{g \inverse},
    \]
    which shows that
    \[
        f \para{g} f\para{g \inverse} = e_H.
    \]

    The result follows from the uniqueness of inverses.
\end{proof}

In other words, \emph{homomorphism are maps between groups that preserve the group structure}.

\begin{definition}[Kernel, Image]
    \label{def:group-homomorphism-kernel-image}%
    Let \(\functionmap{f}{G}{H}\) be a homomorphism. The \emph{kernel} of \(f\) is given by
    \[
        \ker \para{f} = \set{g \in G}{f\para{g} = e_H}
    \]
    and the \emph{image} of \(f\) is given by
    \[
        \img \para{f} = \set{h \in H}{\exists g \in G: f\para{g} = h}.
    \]
\end{definition}

\begin{proposition}[Kernels are Normal Subgroups, Images are Subgroups]
    \label{prop:kernel-normal-subgroup-image-subgroup}%
    Let \(\functiondomain{f}{G}{H}\) be a homomorphism. Then \(\ker\para{f} \normalsub G\) and \(\img\para{f} \subgroup H\).
\end{proposition}

\begin{proof}
    Observe \(e \in \ker\para{f}\) so \(\ker\para{f} \neq \emptyset\). Now, if \(x, y \in \ker\para{f}\), then
    \[
        f\para{x y \inverse} = f\para{x} \cdot f\para{y}\inverse = e \cdot e\inverse = e
    \]
    so \(x y \inverse \in \ker \para{f}\). This shows \(\ker \para{f} \subgroup G\).

    As for the normality, let \(x \in G\) and \(g \in \ker \para{f}\). We calculate
    \[
        f\para{x g x\inverse} = f\para{x} f\para{g} f\para{x}\inverse = f\para{x} f\para{x}\inverse = e
    \]
    so \(x g x\inverse \in \ker \para{f}\). This concludes that \(\ker \para{f} \normalsub G\).

    For the image \(\img \para{f}\), observe that if \(a, b \in \img\para{f}\), say \(a = f\para{x}\) and \(b = f\para{y}\) for \(x, y \in G\), then we have
    \[
        a b \inverse = f\para{x} f\para{y}\inverse = f\para{x y \inverse} \in \img \para{f}
    \]
    and note \(e \in \img \para{f}\), so \(\img \para{f} \subgroup H\) indeed.
\end{proof}

\begin{remark}
    For \(H \normalsub G\), observe that
    \[
        \function{q_H}{G}{G/H}{g}{gH}
    \]
    is a homomorphism.
\end{remark}

\begin{definition}[Group Isomorphism]
    \label{def:group-isomorphism}%
    A group homomorphism \(\functiondomain{f}{G}{H}\) is an \emph{isomorphism} when it is a bijection. Two groups are said to be \emph{isomorphic} if there exists an isomorphism between them, denoted as \(G \isomorphic H\).
\end{definition}

\begin{theorem}[First Isomorphism Theorem]
    \label{thm:first-isomorphism-theorem}%
    Let \(\functiondomain{f}{G}{H}\) be a homomorphism. Then \(\ker \para{f}\) is normal in \(G\), and the homomorphism 
    \[
        \function{\phi}{G / \ker \para{f}}{\img \para{f}}{g \ker \para{f}}{f\para{g}}
    \]
    is an isomorphism. In particular,
    \[
        G / \ker \para{f} \isomorphic \img \para{f}.
    \]
\end{theorem}

\begin{proof}
    We have already shown the normality part previously.

    We first show \(\phi\) is well-defined. Indeed, if \(g \ker f = g' \ker f\), then \(g' g\inverse \in \ker \para{f}\), so \(f\para{g' g\inverse} = e\). This gives \(f\para{g} = f\para{g'}\).

    For the surjectivity of \(\phi\), let \(h \in \img \para{f}\), then \(h = f \para{g}\) for some \(g \in G\), and hence \(h = \phi\para{g \ker f}\).

    For the injectivity of \(\phi\), if \(\phi\para{g \ker f} = \phi\para{g' \ker f}\), then \(f\para{g} = f\para{g'}\), which shows that \(g' g\inverse \in \ker f\). So \(g \ker f = g' \ker f\) and this gives injectivity.
\end{proof}

\begin{example*}
    Consider the homomorphism
    \[
        \function{\exp}{\CC}{\CC\multiplicative}{z}{e^z.}
    \]
    
    The kernel of \(\exp\) is \(2\pi i \ZZ\), so the first isomorphism theorem shows that
    \[
        \CC\multiplicative \isomorphic \CC / 2\pi i \ZZ.
    \]
\end{example*}

\begin{theorem}[Second Isomorphism Theorem]
    \label{thm:second-isomorphism-theorem}%
    Let \(H \subgroup G\) and \(K \normalsub G\). Then
    \[
        HK = \set{hk}{h \in H, k \in K}
    \]
    is a subgroup of \(G\), and \(H \cap K\) is normal in \(H\), then there is an isomorphism
    \[
        HK / K \isomorphic H / H \cap K.
    \]
\end{theorem}

This theorem might seem unnatural at first, but the idea is as follows. We not only want to be able to take the quotient of the original group \(G\) with some normal subgroup \(K \normalsub G\) -- we want to take the quotient with \emph{any} subgroup \(H \subgroup G\), i.e., `\(H / K\)'. But the problem is, \(K\) might not be entirely contained in \(H\). The above two quotients are both choices to make when we take such quotients, and what the theorem states is that they are in fact both isomorphic, to the `quotient'
\[
    \set{hK}{h \in H} \subgroup \set{gK}{g \in G} = G / K.
\]

\begin{proof}
    We first show \(HK\) is a subgroup. It certainly contains \(e = e \cdot e \in HK\) so it is non-empty. Let \(x = hk\) and \(y = h' k'\) for \(x, y \in HK\). We show that \(y x \inverse HK\). Observe that indeed
    \begin{align*}
        y x \inverse &= \para{h' k'}{h k}\inverse \\
        &= h' k' k \inverse h\inverse \\
        &= \para{h' h\inverse} h \para{k' k\inverse} h\inverse \\
        &= \para{h' h\inverse} k'' \\
        &\in HK,
    \end{align*}
    where \(k'' = h \para{k' k\inverse} h\inverse \in K\) since \(K \normalsub G\).

    Consider the homomorphism
    \[
        \function{\phi}{H}{G / H}{h}{hK}
    \]
    by restricting the quotient map. The kernel is
    \[
        \ker \phi = H \cap K.
    \]

    We apply the \autoref{thm:first-isomorphism-theorem} (First Isomorphism Theorem) to \(\phi\), then we have
    \[
        \img \phi \isomorphic H / \ker \phi = H / H \cap K.
    \]

    Now, \(\img \phi\) is the set of \(K\)-cosets in \(G\) that are representable by elements of \(H\), i.e.
    \[
        \set{hK}{h \in H} \subseteq G / K
    \]
    which is precisely \(HK / K\). This finishes the proof.
\end{proof}

\begin{remark}[Correspondence between subgroups of \(G / K\) and \(G\)]
    Let \(\functiondomain{\phi}{G}{G / K}\), the quotient map. There is a bijection between
    \[
        \brce{\text{subgroups of }G/K},
    \]
    and
    \[
        \brce{\text{subgroups of }G\text{ that contain }K}.
    \]

    Given \(K \normalsub L \subgroup G\), we map this to
    \[
        L \mapsto \phi\para{L} = L / K \subgroup G / K.
    \]

    Conversely, let \(A \subgroup G / K\), we map this to
    \[
        \set{g \in G}{gK \in A} = \phi\inverse\para{A} \mapsfrom A.
    \]

    Note that
    \begin{enumerate}
        \item The fact that these are inverses is a tedious but straightforward check;
        \item The correspondence above matches \emph{normal} subgroups on the two sides (again, straightforward check).
    \end{enumerate}
\end{remark}

\begin{theorem}[Third Isomorphism Theorem]
    \label{thm:third-isomorphism-theorem}%
    Let \(K \subgroup L \subgroup G\) be normal subgroups of \(G\). Then there exists an isomorphism
    \[
        \para{G / K} / \para{L / K} \isomorphic G / L.
    \]
\end{theorem}

\begin{proof}
    Define a map
    \[
        \function{\phi}{G / K}{G / L}{gK}{gL}.
    \]

    This map is well-defined. Indeed, if \(gK = g'K\), then we know \(g' g\inverse \in L \subseteq L\), so \(gL = g'L\).

    This is clearly a homomorphism by basic algebra. We now calculate the kernel, which is
    \[
        \ker \phi = \set{gK}{gL = L} = \set{gK}{g \in L} = L / K.
    \]

    This finishes the proof by \autoref{thm:first-isomorphism-theorem}.
\end{proof}

We often study groups by taking quotients of normal subgroups -- but this only works for groups with \emph{non-trivial} normal subgroups.

\begin{definition}[Simplicity]
    \label{def:simplicity}%
    A group \(G\) is called \emph{simple} if the only normal subgroups of \(G\) are \(G\) and \(\brce{e}\).
\end{definition}

\begin{proposition}[Simplicity Condition of Abelian Groups]
    \label{prop:abelian-group-simplicity}%
    Let \(G\) be an abelian group. Then \(G\) is simple if and only if \(G \isomorphic C_p\) for some prime \(p\).
\end{proposition}

\begin{proof}
    If \(G \isomorphic C_p\), then any non-identity element generates \(G\) (by \autoref{thm:lagrange} (Lagrange's Theorem)). So \(G\) is simple.

    Conversely, if \(G\) is simple and abelian, take \(g \in G\) which is not in the identity. We look at the subgroup generated by \(g\), which is
    \[
        \ang{g} = \set{g^n}{n \in \ZZ}
    \]

    This is a subgroup of \(G\), and since \(G\) is abelian, the normality is automatic. Therefore, this subgroup must be \(G\) itself.

    If \(G\) is an infinite cyclic group, then \(G \isomorphic \ZZ\) which is not simple, for example \(2\ZZ \subgroup \ZZ\).

    If \(G\) is finite of order \(n\), say \(G \isomorphic C_m\). If \(q \divides m\) for any \(q\), take the subgroup generated by \(g^{\frac{m}{q}}\), where \(g\) generates \(G\). Normality is then automatic. Therefore, \(q = 1\) or \(m\) by simplicity, so \(m\) is prime.
\end{proof}

\begin{theorem}
    \label{thm:infinite-chain-normal-subgroup}%
    Let \(G\) be a finite group. Then there exists subgroups
    \[
        G = H_1 \normalsup H_2 \normalsup H_3 \normalsup \cdots \normalsup H_n = \brce{e}
    \]
    such that \(H_i / H_{i + 1}\) is simple.
\end{theorem}

\begin{proof}
    If \(G\) is simple, take \(H_2 = \brce{e}\), and we are done.

    Otherwise, let \(H_2\) be a proper normal subgroup of maximal order in \(G\). We claim that \(G / H_2\) is simple. If not, consider the quotient map
    \[
        \functiondomain{\pi}{G}{G / H_2}.
    \]

    Indeed, if \(G / H_2\) is a proper normal subgroup, by the correspondence, we can find such a proper normal subgroup in \(G\) that contains \(H_2\). This contradicts with the maximality of the order of \(\abs{H_2}\).

    We simply repeat the argument to find \(H_3, H_4, \ldots\). Note that \(H_{i + 1} \neq H_i\), and hence \(\abs{H_{i + 1}} \lneq \abs{H_i}\), so the process terminates.
\end{proof}

\begin{remark}
    We also have another infinite class of finite simple groups, that \(A_m\) for \(m \geq 5\) are all simple. We will prove this shortlly.
\end{remark}